% README for the EE101A homework project. This is the project's Main document in Overleaf:
% it is what you see when you press Recompile without opening a homework file.
\documentclass[11pt]{article}
\usepackage[margin=1in]{geometry}
\usepackage{fontspec}
\usepackage{xcolor}
\usepackage{enumitem}
\usepackage{array}
\usepackage{listings}
\usepackage[obeyspaces,spaces]{url}   % before hyperref, which loads url itself
\usepackage[colorlinks=true,linkcolor=blue!60!black,urlcolor=blue!60!black]{hyperref}
\setlist{nosep,leftmargin=1.5em}
\setlength{\parindent}{0pt}
\setlength{\parskip}{0.5em}
\lstset{basicstyle=\ttfamily\small,columns=fullflexible,frame=single,framesep=4pt,
        backgroundcolor=\color{gray!8},rulecolor=\color{gray!40},keepspaces=true,
        aboveskip=0.6em,belowskip=0.6em}
\def\UrlBreaks{\do\/\do\_\do\-}
\let\f\path                                     % file names and keys (line-breakable)
\newcommand{\cmd}[1]{\texttt{\textbackslash\detokenize{#1}}}

\title{EE101A Homework Project --- Start Here}
\author{}
\date{}
\begin{document}
\maketitle
\vspace{-3em}

This project holds every EE101A homework problem (the \emph{question bank}) and the homework
assignments built from them. Problems are written once; each homework lists which problems it uses
and how many points each part is worth. Every PDF comes in a student version and a solutions
version.

\tableofcontents

\section{Compiling a PDF}

\textbf{In Overleaf:} open the file you want in the file tree and press \textbf{Recompile}.
Overleaf compiles the open file, so everyone on the project can compile a different homework at
the same time. (This page is the project's Main document: it is what Recompile shows when the open
file is not one of the files below.)

\begin{center}
\begin{tabular}{>{\raggedright}p{0.30\linewidth} p{0.62\linewidth}}
\hline
\textbf{To get\ldots} & \textbf{open and Recompile} \\ \hline
A homework, student version & \f{assignments/<Semester>/<HWxx>/<Semester>_<HWxx>_student.tex}\newline
                              e.g.\ \f{assignments/2026-4_Fall/HW07/2026-4_Fall_HW07_student.tex} \\
A homework, solutions & the \f{..._solutions.tex} file next to it \\
Every problem in the bank & \f{main_questionbank_student.tex} or \f{main_questionbank_solutions.tex} \\
\hline
\end{tabular}
\end{center}

The compiler must be \textbf{XeLaTeX} (Menu $\rightarrow$ Compiler). The first compile of a
homework can show 0 points; compiling again fixes it (Overleaf usually does this automatically).

\section{Where things are}

\begin{lstlisting}
assignments/<Semester>/          one folder per offering, sorted in time order:
                                 <year>-<n>_<Season>; n: 1 Winter ... 4 Fall
                                 e.g. 2026-1_Winter, 2026-4_Fall
    course.tex                   printed name ("Fall 2026"), instructor, style
    HW07/
        metadata.tex             homework number and due date
        questions.tex            which problems, in what order, with what points
        2026-4_Fall_HW07_student.tex / ..._solutions.tex   <- compile these
questionbank/
    04_diode_circuits/           one folder per syllabus topic (numbered in course order)
        load_line/               subtopic: the main concept a problem teaches
            led_iv_curve/        one problem; this folder name is its "id"
                v0/ v1/          versions of the problem (source.tex + figures)
                version_history.md   author, date, notes, concept tags
    README.md                    table of every problem (generated)
    tags.md                      the list of concept tags
config/  style/                  shared setup and layout -- rarely edited
tools/                           scripts (section 5)
\end{lstlisting}

\section{Common tasks}

\subsection{Put an existing problem in a homework}
Find it in the question bank PDF (\f{main_questionbank_solutions.tex}): each version's box shows
its \emph{id}, \emph{version} and \emph{part keys}. Add a block to the homework's
\f{questions.tex}:
\begin{lstlisting}
% Problem 3 - LED I-V Curve Load-Line Analysis (15 points)
\clearpage
\UseQuestion{led_iv_curve}{v1}{
    basicloadline     = 3,
    ledbrightness     = 2,
    resistorselection = 5,
    ambienthot        = 5 extra,     % "extra" = extra credit
}
\end{lstlisting}
Every part must get points. If a problem has no parts, or its points are only given as a total,
write \f{total = 10,} instead of the part list.

\subsection{Add a new problem}
\begin{enumerate}
  \item Pick the subtopic folder for the concept the problem mainly teaches, and create
        \f{<id>/v0/source.tex} with a short, descriptive snake\_case id (e.g.\
        \f{questionbank/04_diode_circuits/load_line/zener_load_line/v0/source.tex}). Put its
        figures in the same \f{v0/} folder.
  \item Write the problem in this shape:
\begin{lstlisting}
\begin{mainquestion}{zenerloadline}{Zener Load Line}
Problem statement, given values, figures...
\begin{center}\includegraphics[width=3in]{circuit.png}\end{center}

\begin{partquestion}{findcurrent}
{Prompt for part (a). Letters and points are added automatically.}
\begin{answer}[2in]
Solution for part (a).
\end{answer}
\end{partquestion}
\end{mainquestion}
\end{lstlisting}
        Keep \cmd{begin{answer}} and \cmd{end{answer}} on their own lines with no indentation;
        figures are referred to by file name only.
  \item Add \f{<id>/version_history.md} (copy one from a neighbouring problem): title, a
        \f{Concepts:} line using tags from \f{questionbank/tags.md}, and a \texttt{\#\#\# v0} entry with
        author, date and a one-line description.
  \item Use it in a homework (section 3.1), then \textbf{update the catalog}: run
        \f{python tools/check.py --fix-index --catalog} and \f{python tools/check.py}
        (section 5). Without this step the new problem is not found and does not appear in the
        question bank PDF.
\end{enumerate}

\subsection{Change an existing problem}
If students have already received it, do not edit it: copy \f{v0/} to \f{v1/}, change the copy,
add a \texttt{\#\#\# v1} entry to \f{version_history.md}, and point the new homework at \f{v1}. Typos
may be fixed in place.

\subsection{Create a new homework or semester}
Copy an existing \f{HWxx/} folder, edit \f{metadata.tex} (number, due date) and
\f{questions.tex}, then run \f{python tools/check.py --fix-mains} to create the two files you
compile. For a new semester, create its folder with the dated name (e.g.\ \f{2027-1_Winter}),
copy a \f{course.tex} into it, and set \cmd{SemesterName} (e.g.\ \f{Winter 2027}) and the
instructor.

\subsection{Change the look}
In a semester's \f{course.tex}:
\begin{itemize}
  \item \cmd{HomeworkStyle{style-jeff}} --- the standard layout;
  \item \cmd{HomeworkStyle{style-mark}} --- Mark Horowitz's Winter 2026 layout;
  \item \cmd{boxedsolutionsfalse} --- removes the blue frame around style-jeff solutions.
\end{itemize}
A single homework can override these in its \f{metadata.tex}.

\subsection{Move or rename a problem}
Move its folder, then run \f{python tools/check.py --fix-index --catalog}. Homeworks refer to
problems by id, so a move needs nothing else; a rename also needs the id changed in every
\f{questions.tex} that uses it.

\section{Quick answers}
\begin{description}[style=nextline,leftmargin=1.5em]
  \item[A part shows 0 points or a total looks wrong.] Compile again (totals need two passes);
        if it persists, a part key in \f{questions.tex} doesn't match the problem's
        \cmd{begin{partquestion}{...}} --- \f{tools/check.py} lists every mismatch.
  \item[``Unknown question id''.] The id isn't in \f{questionbank/index.tex}: run
        \f{python tools/check.py --fix-index}.
  \item[Where did a problem get used before?] The question bank PDF and
        \f{questionbank/README.md} list every homework that used each version.
  \item[Can I break a past homework?] Only by editing a version it uses. Make a new version
        instead (section 3.3).
  \item[More detail?] \f{import/import_instructions_for_LLM.md} documents every convention;
        it is also what an AI assistant follows when importing new problems.
\end{description}

\section{Scripted tools}

The scripts run on a local copy of the project (git checkout with Python), not inside Overleaf.
After running them, upload the changed files to Overleaf.

\begin{center}
\begin{tabular}{>{\raggedright}p{0.36\linewidth} p{0.56\linewidth}}
\hline
\textbf{Command} & \textbf{When to use it} \\ \hline
\f{python tools/check.py} & \textbf{After any change}, before uploading. Checks points vs.\
    parts, ids, figures, cross-references, version histories and tags without compiling.
    Must report 0 errors. \\
\f{python tools/check.py --fix-index} & After adding, moving or renaming a problem folder
    (rewrites \f{questionbank/index.tex}). \\
\f{python tools/check.py --catalog} & After adding problems, editing \f{version_history.md}
    or tags, or changing a homework's problems: refreshes the question bank PDF contents and
    \f{questionbank/README.md}. \\
\f{python tools/check.py --fix-mains} & After creating a homework folder: writes its
    \f{_student.tex} / \f{_solutions.tex} files. \\
\f{python tools/check.py --rename-tag old new} & Renames a concept tag everywhere. \\
\texttt{.\textbackslash build.ps1 -Semester 2026-4\_Fall -Assignment HW07 [-Mode student]} & Compiles one
    homework locally to \f{out.pdf} (Windows; one-time setup \texttt{.\textbackslash build.ps1 -Install}). In VS
    Code: open a \f{_student/_solutions.tex} file and press Ctrl+Shift+B. \\
\f{python tools/build_all.py} & Compiles every homework and the question bank. Use after
    changing shared files (\f{style/}, \f{config/}); \f{--only 2026-4_Fall/HW07} for one
    homework, \f{--save}/\f{--compare} to check nothing else changed. \\
\hline
\end{tabular}
\end{center}

\end{document}
